\chapter*{Abstract}
\addcontentsline{toc}{chapter}{Abstract}

This textbook presents \emph{DAGonStar} (Direct Acyclic Graph On anything, also
written \emph{DAGon*}), a lightweight Python workflow engine designed to
orchestrate computational tasks across local machines, on-premise
high-performance computing (HPC) clusters, containers, cloud-based virtual
infrastructures, and, in its current evolution, artificial-intelligence (AI)
and web services. The material is organized as a self-contained academic
treatment of workflow-based computational science, structured around the same
progressive narrative used by the project's own fourteen-lesson tutorial
sequence, and grounded exclusively in the documentation and source examples
found in the \texttt{docs/} directory of the
\texttt{DagOnStar/dagonstar} repository
(\texttt{development-native-web} branch).

The book is organized in three parts. Part~I establishes the theoretical
foundations of scientific workflow orchestration --- directed acyclic graphs,
task- and data-oriented dependency models, the \texttt{workflow://} schema,
configuration management, and the operational mechanics of scratch
directories, generated launchers, cycle detection, checkpointing, and data
staging. Part~II extends the model to heterogeneous execution backends:
containers, remote and Slurm-scheduled HPC execution, meta-workflows,
asynchronous orchestration, and --- distinctively --- first-class
\emph{AI--HPC} integration through Large Language Model (LLM) tasks, Web
tasks, and Native Python tasks. Part~III applies the accumulated theory to
computational environmental science, using the project's WRF/ROMS/WaComM++-inspired
environmental application as an extended case study, and closes with a
synthesis chapter on operational, reproducible, AI-augmented forecasting
chains. A set of appendices collects the complete API reference, developer
practices, configuration reference, and checkpoint semantics for use as a
standing reference after the narrative chapters have been read.

Every code listing, configuration fragment, and lesson narrative reproduced in
this book is drawn directly from the DAGonStar documentation and tutorial
sources; the book adds pedagogical structure, formal definitions, discussion,
and original explanatory figures, while remaining descriptive of, and
consistent with, the current implementation.
